\documentclass[aspectratio=169,11pt]{beamer}

\usepackage{iftex}
\ifPDFTeX
  \usepackage[utf8]{inputenc}
  \usepackage[T5]{fontenc}
  \usepackage[vietnamese]{babel}
\else
  \usepackage{fontspec}
  \setsansfont{TeX Gyre Heros}
  \setmainfont{TeX Gyre Termes}
  \setmonofont{Inconsolata}
  \usepackage[vietnamese]{babel}
\fi

\usepackage{graphicx}
\usepackage{booktabs}
\usepackage{tikz}
\usepackage{listings}
\usepackage{xcolor}
\usepackage{hyperref}
\usetikzlibrary{shapes.geometric, arrows.meta, positioning, shadows}
\usepackage{pifont}


\vfuzz=120pt
\hbadness=10000
\hfuzz=2pt
\pdfstringdefDisableCommands{%
  \def\textbf#1{#1}%
  \def\\{ }%
}
\DeclareFontShape{T5}{cmss}{m}{it}{<->ssub*cmss/m/sl}{}

\definecolor{hcmusblue}{HTML}{005BAA}
\definecolor{hcmusred}{HTML}{D71920}
\definecolor{darktext}{HTML}{202428}
\definecolor{primaryblue}{RGB}{41, 128, 185}
\definecolor{softgray}{RGB}{236, 240, 241}
\definecolor{accentorange}{RGB}{230, 126, 34}
\definecolor{successgreen}{RGB}{39, 174, 96}
\definecolor{dangerred}{RGB}{192, 57, 43}

\usetheme{Madrid}
\usecolortheme{default}
\setbeamercolor{structure}{fg=hcmusblue}
\setbeamercolor{title}{fg=white,bg=hcmusblue}
\setbeamercolor{frametitle}{fg=white,bg=hcmusblue}
\setbeamercolor{block title}{fg=white,bg=hcmusblue}
\setbeamercolor{block body}{bg=softgray,fg=darktext}
\setbeamertemplate{navigation symbols}{}
\setbeamertemplate{footline}[frame number]

\newcommand{\ProjectTitle}{Hệ thống phát hiện đạo văn}
\newcommand{\CourseName}{Khai thác dữ liệu văn bản và ứng dụng}
\newcommand{\HCMUSLogo}{%
  \IfFileExists{figures/hcmus-logo.png}{%
    \includegraphics[height=0.9cm]{figures/hcmus-logo.png}%
  }{%
    \begin{tikzpicture}[scale=0.36]
      \draw[line width=1.2pt,color=hcmusblue] (0,0) circle (1.25);
      \draw[line width=0.8pt,color=hcmusred] (0,0) circle (0.95);
      \node[font=\bfseries\scriptsize,color=hcmusblue] at (0,0) {HCMUS};
    \end{tikzpicture}%
  }%
}

\title[\ProjectTitle]{\textbf{\ProjectTitle}}
\subtitle{\CourseName}

\author[Nhóm 3]{\texorpdfstring{%
  \small\textbf{GVHD: Ph.D Lê Thanh Tùng}\\[0.25cm]
  \textbf{Nhóm 3}\\[0.1cm]
  \scriptsize
  \begin{tabular}{l}
    \textbf{Sinh viên thực hiện:} \\
    23127004 - Lê Nhật Khôi \\
    23127049 - Đỗ Hoàng Duy Hưng \\
    23127189 - Trần Trọng Hiếu \\
    23127210 - Nguyễn Đăng Khoa \\
    23127248 - Nguyễn Hữu Phúc
  \end{tabular}
}{Nhóm 3}}
\hypersetup{pdfauthor={Nhóm 3},pdftitle={Hệ thống phát hiện đạo văn}}

\institute[HCMUS]{
  Khoa Công nghệ thông tin\\
  Trường Đại học Khoa học tự nhiên, ĐHQG-HCM
}
\date{Tháng 07 năm 2026}

\begin{document}

\begin{frame}
  \titlepage
  \begin{tikzpicture}[remember picture,overlay]
    \node[anchor=north east,xshift=-0.55cm,yshift=-0.85cm] at (current page.north east) {\HCMUSLogo};
  \end{tikzpicture}
\end{frame}

\begin{frame}{Nội dung}
\tableofcontents
\end{frame}

\section{Bài toán}
{
\setbeamercolor{background canvas}{bg=white} % Đặt nền trắng

\begin{frame}[plain]
    \vfill
    \centering
    % Dùng lệnh \textcolor để ép màu xanh hcmusblue cho dòng chữ
    \textcolor{hcmusblue}{\textbf{\Huge Bài toán}}
    \vfill
\end{frame}
}

\begin{frame}{Bối cảnh bài toán}
\setbeamercolor{block title}{bg=dangerred,fg=white}
\begin{block}{Vấn đề thực tiễn}
Trong kỷ nguyên số, tài liệu học thuật và nội dung trực tuyến có thể bị sao chép, diễn giải lại hoặc được LLM viết lại rất khó phát hiện bằng so khớp từ khóa thông thường.
\end{block}

\vspace{0.25cm}
\begin{columns}[T]
\begin{column}{0.48\textwidth}
\setbeamercolor{block title}{bg=accentorange,fg=white}
\begin{block}{Hai dạng đạo văn chính}
\begin{itemize}
  \item[\ding{55}] \textbf{Đạo văn truyền thống:} sao chép nguyên văn đoạn văn, số liệu hoặc ý tưởng nhưng không trích dẫn.
  \item[\ding{55}] \textbf{Đạo văn tạo sinh:} dùng LLM paraphrase lại cách diễn đạt, nhưng vẫn giữ cấu trúc logic và ý tưởng cốt lõi.
\end{itemize}
\end{block}
\end{column}

\begin{column}{0.48\textwidth}
\setbeamercolor{block title}{bg=successgreen,fg=white}
\begin{block}{Nhu cầu của hệ thống}
\begin{itemize}
  \item[\ding{51}] Phát hiện không chỉ trùng lặp bề mặt mà cả \textbf{tương đồng ngữ nghĩa}.
  \item[\ding{51}] Phân biệt văn bản AI viết độc lập với văn bản AI viết dựa trên nguồn có sẵn.
  \item[\ding{51}] Chỉ ra \textbf{đạo bao nhiêu} và \textbf{đạo từ nguồn nào} để đối chiếu.
\end{itemize}
\end{block}
\end{column}
\end{columns}
\end{frame}

\begin{frame}{Phát biểu bài toán}
\textbf{\Large \textcolor{primaryblue}{Phát hiện đạo văn trên tài liệu đầu vào}}
\vspace{0.15cm}

\setbeamercolor{block title}{bg=primaryblue,fg=white}
\begin{block}{Input -- Output}
\begin{columns}[T]
\begin{column}{0.48\textwidth}
\textbf{Input}
\begin{itemize}
  \item Một tài liệu nghi vấn $P$.
  \item Kho tài liệu nguồn $\mathcal{S}$ để truy vết nguồn đạo văn.
\end{itemize}
\end{column}
\begin{column}{0.48\textwidth}
\textbf{Output}
\begin{itemize}
  \item Điểm đạo văn trong khoảng \textbf{0--1}.
  \item Tỷ lệ phần trăm nội dung bị nghi ngờ đạo văn.
  \item Danh sách nguồn và đoạn văn tương ứng.
\end{itemize}
\end{column}
\end{columns}
\end{block}

\vspace{0.15cm}
\setbeamercolor{block title}{bg=accentorange,fg=white}
\begin{block}{Dạng kết quả cần xác định}
\small
Với tài liệu nghi vấn $P$ và kho nguồn $\mathcal{S}$, hệ thống cần tìm tập các trường hợp đạo văn:
\[
  \mathcal{D}=\{d=(S, s, p)\mid S\in\mathcal{S},\ s\in S,\ p\in P\}
\]
Trong đó $S$ là tài liệu nguồn, $s$ là đoạn gốc và $p$ là đoạn nghi vấn trong tài liệu đầu vào.
\end{block}
\end{frame}

\begin{frame}{Quy trình giải quyết bài toán}
\begin{columns}[T]
\begin{column}{0.32\textwidth}
\setbeamercolor{block title}{bg=primaryblue,fg=white}
\begin{block}{1. Source Retrieval}
\footnotesize
Sàng lọc từ kho tri thức để tìm các tài liệu nguồn nghi vấn có liên quan đến tài liệu đầu vào $P$.
\end{block}
\end{column}

\begin{column}{0.32\textwidth}
\setbeamercolor{block title}{bg=accentorange,fg=white}
\begin{block}{2. Text Alignment}
\footnotesize
So khớp từng cặp $(S,P)$ để xác định đoạn nguồn $s$ và đoạn nghi vấn $p$ bằng offset và độ dài.
\end{block}
\end{column}

\begin{column}{0.32\textwidth}
\setbeamercolor{block title}{bg=successgreen,fg=white}
\begin{block}{3. Scoring}
\footnotesize
Tổng hợp các đoạn phát hiện được để chấm điểm 0--1, tính tỷ lệ đạo văn và chỉ ra nguồn tương ứng.
\end{block}
\end{column}
\end{columns}

\vspace{0.25cm}
\setbeamercolor{block title}{bg=dangerred,fg=white}
\begin{block}{Minh họa đầu ra}
\small
\textbf{Tài liệu đầu vào} $\rightarrow$ \textbf{Điểm đạo văn:} 0.72 \; $\rightarrow$ \; \textbf{Tỷ lệ:} 72\% nội dung nghi vấn \; $\rightarrow$ \; \textbf{Nguồn:} \texttt{source-document-ABC.txt}, vị trí đoạn nguồn và đoạn nghi vấn.
\end{block}
\end{frame}
\begin{frame}{Kiến trúc data pipeline}
\begin{figure}[htbp]
\centering
\begin{tikzpicture}[
  node distance=1.55cm,
  box/.style={draw=primaryblue!80, fill=primaryblue!10, thick, rounded corners=4pt, align=center, minimum width=6cm, minimum height=0.75cm, font=\small\bfseries},
  io/.style={draw=accentorange!80, fill=accentorange!10, thick, rounded corners=4pt, align=center, minimum width=6cm, minimum height=0.7cm, font=\small},
  gate/.style={draw=accentorange!90, fill=accentorange!12, very thick, rounded corners=4pt, align=center, minimum width=6cm, minimum height=0.75cm, font=\small\bfseries},
  future/.style={draw=gray!60, fill=gray!8, thick, dashed, rounded corners=4pt, align=center, minimum width=6cm, minimum height=0.7cm, font=\small},
  arrow/.style={-Latex, line width=1pt, primaryblue},
  lbl/.style={right, font=\scriptsize, text=black!70}
]
% Nodes
\node[io] (raw) {Raw PAN 2025 dataset\\{\footnotesize susp/ + src/ + pairs + *\_truth XML}};
\node[gate, below of=raw] (integ) {Integrity Check (A0-2)\\{\footnotesize scripts/check\_integrity.py}};
\node[io, below of=integ] (valid) {Validated splits\\{\footnotesize documents + ground-truth XML}};
\node[box, below of=valid] (parser) {Label Parser (A0-3)\\{\footnotesize scripts/parse\_labels.py}};
\node[io, below of=parser, fill=successgreen!10, draw=successgreen!80] (out) {Structured labels\\{\footnotesize *\_labels.jsonl + *\_spans.csv}};
\node[future, below of=out] (down) {Downstream --- Phase 1--3\\{\footnotesize EDA $\cdot$ Retrieval $\cdot$ Alignment $\cdot$ PlagDet}};

% Arrows
\draw[arrow] (raw) -- node[lbl] {3 splits} (integ);
\draw[arrow] (integ) -- node[lbl] {PASS: 0 missing / 0 orphan} (valid);
\draw[arrow] (valid) -- node[lbl] {XML only} (parser);
\draw[arrow] (parser) -- node[lbl] {2{,}730{,}100 spans} (out);
\draw[arrow, dashed, gray] (out) -- node[lbl] {upcoming} (down);
\end{tikzpicture}
\caption{Luồng dữ liệu Phase 0: bộ dữ liệu thô PAN 2025 được kiểm tra tính toàn vẹn (A0-2), sau đó nhãn XML được phân tích thành span có cấu trúc (A0-3) phục vụ các pha truy hồi, căn chỉnh và chấm điểm phía sau.}
\label{fig:data-pipeline}
\end{figure}
\end{frame}
% \section{Dataset tổng kết}
% {
% \setbeamercolor{background canvas}{bg=white} % Đặt nền trắng

% \begin{frame}[plain]
%     \vfill
%     \centering
%     % Dùng lệnh \textcolor để ép màu xanh hcmusblue cho dòng chữ
%     \textcolor{hcmusblue}{\textbf{\Huge Dataset tổng kết}}
%     \vfill
% \end{frame}
% }

% \begin{frame}{Dataset Luật và Câu hỏi}
%     % Phần mô tả nằm trên, chiếm hết bề ngang
%     \setbeamercolor{block title}{bg=primaryblue,fg=white}
%     \begin{block}{Mô tả tập dữ liệu Văn bản luật}
%         \small
%         \begin{itemize}
%             \item \textbf{Nguồn:} Bộ luật Lao động VN (Số: 125/VBHN-VPQH, 2025). \textbf{Quy mô:} \textbf{220} điều, \textbf{17} chương.
%             \item \textbf{Cấu trúc RAG:} Phân rã thành các chunks, tự động gán nhãn các chủ đề pháp lý khác nhau.
%             \item \textbf{Metadata:} Lưu trữ chi tiết về chương, tên chương, tên điều và chủ đề để hỗ trợ lọc thông minh.
%         \end{itemize}
%     \end{block}

%     \vspace{0.1cm}
%     \setbeamercolor{block title}{bg=accentorange,fg=white}
%     \begin{block}{Mô tả tập dữ liệu Câu hỏi}
%         \small
%         \begin{itemize}
%             \item \textbf{Số lượng:} \textbf{4.899} câu hỏi đã làm sạch.
%             \item \textbf{Mapping:} Mỗi câu map với \texttt{article\_id} làm Ground Truth cho việc truy xuất.
%         \end{itemize}
%     \end{block}
% \end{frame}

% \begin{frame}{Quy trình chuẩn hoá dữ liệu}
% \setbeamercolor{block title}{bg=dangerred,fg=white}
% \begin{block}{Vấn đề ban đầu}
% \begin{itemize}
%   \item[\ding{55}] \texttt{law.json}: Các điều luật có "câu dẫn" bị lệch khoản khi extract trực tiếp sang JSON.
%   \item[\ding{55}] \texttt{questions.json}: Tồn tại \textbf{19} câu hỏi sai định dạng cấu trúc trong tổng 4.918 câu gốc.
% \end{itemize}
% \end{block}

% \vspace{0.2cm}
% \setbeamercolor{block title}{bg=successgreen,fg=white}
% \begin{block}{Giải pháp xử lý (Cleaning Pipeline)}
% \begin{itemize}
%   \item[\ding{51}] Viết script sửa lại logic parsing để \textbf{giữ nguyên cấu trúc JSON chuẩn} của các điều có câu dẫn.
%   \item[\ding{51}] \textbf{Loại bỏ toàn bộ câu sai định dạng} $\rightarrow$ \textbf{4.899 câu hỏi}.
%   \item[\ding{51}] Gán lại nhãn \texttt{article\_id} mới, lấy ngẫu nhiên \textbf{250 mẫu đối chiếu thủ công} để kiểm tra độ chính xác.
% \end{itemize}
% \end{block}
% \end{frame}

% \begin{frame}{Minh họa cấu trúc dữ liệu thực tế}
%     \centering
%     \setbeamercolor{block title}{bg=softgray,fg=primaryblue}
%     \begin{block}{\footnotesize Cấu trúc khoản luật}
%         \centering % Lệnh này căn giữa nội dung bên trong block
%         \includegraphics[width=0.95\textwidth, height=0.38\textheight, keepaspectratio]{figures/law_example.png}
%     \end{block}

%     \vspace{0.1cm}

%     \centering
%     \setbeamercolor{block title}{bg=softgray,fg=accentorange}
%     \begin{block}{\footnotesize Cấu trúc một câu hỏi}
%         \centering
%         \includegraphics[width=0.95\textwidth, height=0.38\textheight, keepaspectratio]{figures/question_example.png}
%     \end{block}
% \end{frame}

% \begin{frame}{Quy trình xây dựng tập huấn luyện}
% \setbeamercolor{block title}{bg=primaryblue,fg=white}
% \begin{block}{Quy trình chuyển đổi định dạng (Formatting)}
% \begin{enumerate}
%     \item \textbf{Input:} \texttt{questions.json} chứa dữ liệu thô (Câu hỏi - Ground Truth).
%     \item \textbf{Data Enrichment:} Sử dụng RAG Engine để truy xuất văn bản luật liên quan.
%     \item \textbf{Conversion:} Áp dụng template ChatML để đóng gói dữ liệu theo định dạng [System - User - Assistant].
%     \item \textbf{Split:} Phân chia tập dữ liệu thành Train/Val/Test.
% \end{enumerate}
% \end{block}

% \vspace{0.2cm}
% \centering
% \setbeamercolor{block title}{bg=softgray,fg=primaryblue}
% \begin{block}{\footnotesize Cấu trúc dữ liệu câu hỏi sau chuyển đổi (ChatML)}
%     \centering
%     \includegraphics[width=0.9\textwidth, height=0.35\textheight, keepaspectratio]{figures/chatml_example.png}
% \end{block}
% \end{frame}

% \begin{frame}{Phân chia tập dữ liệu \& Cấu hình Fine-tuning LLM}
% \begin{columns}[T]
% \begin{column}{0.55\textwidth}
% \setbeamercolor{block title}{bg=primaryblue,fg=white}
% \begin{block}{Split Dataset (Tổng: 4.899 câu)}
% \centering
% \renewcommand{\arraystretch}{1.3}
% \begin{tabular}{lrr}
% \toprule
% \textbf{Tập dữ liệu} & \textbf{Số lượng} & \textbf{Tỉ lệ}\\
% \midrule
% \textbf{Train} & 4.400 câu & 89,8\%\\
% \textbf{Validation} & 250 câu & 5,1\%\\
% \textbf{Test} & 249 câu & 5,1\%\\
% \bottomrule
% \end{tabular}
% \end{block}
% \end{column}

% \begin{column}{0.43\textwidth}
% \setbeamercolor{block title}{bg=accentorange,fg=white}
% \begin{block}{Hệ thống Fine-tuning}
% \begin{itemize}
%   \item \textbf{Mô hình:} Qwen2.5-1.5B (nhẹ \& mạnh xử lý tiếng Việt)
%   \item \textbf{Phương pháp:} LoRA (via Unsloth)
%   \item \textbf{Format:} ChatML
%   \item \textbf{Lưu trữ:} Lưu checkpoints ở 275, 550, 825 steps.
% \end{itemize}
% \end{block}
% \end{column}
% \end{columns}
% \end{frame}


\section{Dataset tổng kết}
{
\setbeamercolor{background canvas}{bg=white} % Đặt nền trắng

\begin{frame}[plain]
    \vfill
    \centering
    % Dùng lệnh \textcolor để ép màu xanh hcmusblue cho dòng chữ
    \textcolor{hcmusblue}{\textbf{\Huge Dataset tổng kết}}
    \vfill
\end{frame}
}

\begin{frame}{Giới thiệu Dataset -- PAN 2025}
\setbeamercolor{block title}{bg=primaryblue,fg=white}
\begin{block}{Nguồn dữ liệu}
\small
\textbf{PAN 2025 -- Generated Plagiarism Detection} (Task A0-2): các cặp tài liệu \textbf{suspicious / source} với nhãn đạo văn được sinh bởi LLM (DeepSeek-R1, Llama-3, Mistral) ở ba mức obfuscation \texttt{simple}/\texttt{medium}/\texttt{hard}.
\end{block}

\vspace{0.2cm}
\begin{columns}[T]
\begin{column}{0.48\textwidth}
\setbeamercolor{block title}{bg=accentorange,fg=white}
\begin{block}{Quy mô dữ liệu}
\begin{itemize}
  \item[\ding{51}] \textbf{Train:} 60.759 susp, 60.592 src, 62.160 pairs.
  \item[\ding{51}] \textbf{Validation:} 7.950 susp, 7.949 src, 7.976 pairs.
  \item[\ding{51}] \textbf{Spot-check:} 50 cặp để kiểm tra nhanh (sanity check).
\end{itemize}
\end{block}
\end{column}

\begin{column}{0.48\textwidth}
\setbeamercolor{block title}{bg=successgreen,fg=white}
\begin{block}{Đặc điểm nhãn}
\begin{itemize}
  \item[\ding{51}] \textbf{plagiarism:} đoạn có nguồn tương ứng (\texttt{source\_reference}).
  \item[\ding{51}] \textbf{altered:} đoạn do LLM viết lại, không có nguồn.
  \item[\ding{51}] \textbf{original:} Không có đạo văn.
  \item[\ding{51}] Mức \textbf{severity} tài liệu: \texttt{low}/\texttt{medium}/\texttt{high}.
\end{itemize}
\end{block}
\end{column}
\end{columns}
\end{frame}

\begin{frame}[fragile]{Cấu trúc thư mục \& vai trò tệp tin}
\setbeamercolor{block title}{bg=primaryblue,fg=white}
\begin{block}{Cấu trúc thư mục}
\scriptsize
\begin{verbatim}
*/
|-- */ 
|   |-- susp/suspicious-documentXXXXXX.txt
|   |-- src/source-documentXXXXXX.txt
|   \-- pairs
\-- *_truth/
    |-- metadata.json
    \-- *.xml
\end{verbatim}
\end{block}

\vspace{0.1cm}
\setbeamercolor{block title}{bg=softgray,fg=darktext}
\begin{block}{\footnotesize Vai trò các tệp chính}
\scriptsize
\renewcommand{\arraystretch}{1.2}
\begin{tabular}{ll}
\toprule
\textbf{Tệp / Thư mục} & \textbf{Nội dung} \\
\midrule
\texttt{susp/suspicious-documentNNNNNN.txt} & Tài liệu nghi vấn (UTF-8) \\
\texttt{src/source-documentNNNNNN.txt} & Tài liệu nguồn khả nghi \\
\texttt{pairs} & Danh sách cặp \texttt{<susp> <src>} ứng viên \\
\texttt{*\_truth/*.xml} & Nhãn ground-truth cho từng cặp \\
\texttt{*\_truth/metadata.json} & Metadata arXiv \\
\bottomrule
\end{tabular}
\end{block}
\end{frame}

\begin{frame}{Schema nhãn Ground-truth (XML)}
\begin{block}{}
\begin{columns}[T]
\begin{column}{0.48\textwidth}
\centering
\includegraphics[height=3.8cm,keepaspectratio]{figures/xml.png}
\end{column}
\begin{column}{0.48\textwidth}
\centering
\includegraphics[height=3.8cm,keepaspectratio]{figures/metadata.jpg}
\end{column}
\end{columns}
\end{block}

\vspace{0.07cm}
\setbeamercolor{block title}{bg=accentorange,fg=white}
\begin{block}{Các \texttt{feature} trong một tệp XML}
\scriptsize
\renewcommand{\arraystretch}{1.25}
\begin{tabular}{lll}
\toprule
\textbf{Feature} & \textbf{Ý nghĩa} & \textbf{Thuộc tính chính} \\
\midrule
\texttt{about} & Metadata cấp tài liệu & \texttt{similarity}, \texttt{severity} \\
\texttt{md5Hash} & Hash toàn vẹn tài liệu nghi vấn & \texttt{value} \\
\textbf{\texttt{plagiarism}} & Đoạn đạo văn \textbf{có nguồn} & \texttt{this\_offset/length}, \texttt{source\_reference}, \texttt{obfuscation} \\
\textbf{\texttt{altered}} & Đoạn LLM viết lại, \textbf{không nguồn} & \texttt{this\_offset/length} \\
\bottomrule
\end{tabular}
\end{block}
\end{frame}

% \begin{frame}{Quy mô dataset \& Kiểm tra toàn vẹn}
% \begin{columns}[T]
% \begin{column}{0.55\textwidth}
% \setbeamercolor{block title}{bg=primaryblue,fg=white}
% \begin{block}{Thống kê theo split}
% \centering
% \scriptsize
% \renewcommand{\arraystretch}{1.2}
% \begin{tabular}{lrrr}
% \toprule
% \textbf{Split} & \textbf{Pairs} & \textbf{Truth XML} & \textbf{Spans} \\
% \midrule
% Spot-check & 50 & 50 & 1.639 \\
% Train & 62.160 & 62.160 & 2.730.100 \\
% Validation & 7.976 & 348.629 & 348.629 \\
% \bottomrule
% \end{tabular}

% \vspace{0.15cm}
% \footnotesize Train: 1.877.750 plagiarism / 852.350 altered \\

% Validation: 238.242 plagiarism / 110.387 altered
% \end{block}
% \end{column}

% \begin{column}{0.43\textwidth}
% \setbeamercolor{block title}{bg=successgreen,fg=white}
% \begin{block}{Integrity check}
% \scriptsize
% \begin{itemize}
%   \item[\ding{51}] 0 file thiếu, 0 file mồ côi (orphan).
%   \item[\ding{51}] số \texttt{pairs} bằng số XML ở mọi split.
%   \item[\ding{51}] Susp/src tái sử dụng: train 1.401/1.568; validation 26/27.
%   \item[\ding{51}] Không có \texttt{.txt} nào bị bỏ sót khỏi \texttt{pairs}.
% \end{itemize}
% \end{block}
% \end{column}
% \end{columns}
% \end{frame}

\begin{frame}{Phân bố nhãn \& Obfuscation (tập Train)}
\begin{columns}[T]
\begin{column}{0.5\textwidth}
\setbeamercolor{block title}{bg=primaryblue,fg=white}
\begin{block}{Loại đoạn \& mức độ nghiêm trọng}
\centering
\scriptsize
\renewcommand{\arraystretch}{1.2}
\begin{tabular}{lr}
\toprule
\textbf{Loại đoạn (pair)} & \textbf{Tỷ lệ} \\
\midrule
Plagiarism & 69,4\% \\
Altered & 24,3\% \\
Original & 6,3\% \\
\midrule
\textbf{Severity} & \textbf{Tỷ lệ} \\
\midrule
Medium & 37,4\% \\
Low & 31,4\% \\
High & 24,9\% \\
\bottomrule
\end{tabular}
\end{block}
\end{column}

\begin{column}{0.45\textwidth}
\setbeamercolor{block title}{bg=accentorange,fg=white}
\begin{block}{Obfuscation \& LLM sinh văn bản}
\centering
\scriptsize
\renewcommand{\arraystretch}{1.2}
\begin{tabular}{lr}
\toprule
\textbf{Obfuscation} & \textbf{Tỷ lệ} \\
\midrule
Simple & 60,0\% \\
Medium & 30,0\% \\
Hard & 10,0\% \\
\midrule
\textbf{LLM} & \textbf{Tỷ lệ} \\
\midrule
Llama-3 & 44,4\% \\
DeepSeek-R1 & 42,5\% \\
Mistral & 13,1\% \\
\bottomrule
\end{tabular}
\end{block}
\end{column}
\end{columns}
\end{frame}

\begin{frame}{EDA -- Độ dài đoạn \& số đoạn mỗi cặp (tập Train)}
\setbeamercolor{block title}{bg=softgray,fg=darktext}
\begin{block}{\footnotesize Thống kê độ dài đoạn (ký tự)}
\centering
\scriptsize
\renewcommand{\arraystretch}{1.2}
\begin{tabular}{lrrrrr}
\toprule
\textbf{Loại đoạn} & \textbf{Mean} & \textbf{Median} & \textbf{P25} & \textbf{P75} & \textbf{P90} \\
\midrule
Plagiarism & 785 & 685 & 395 & 1.021 & 1.425 \\
Altered & 714 & 583 & 340 & 926 & 1.348 \\
\bottomrule
\end{tabular}
\end{block}

\vspace{0.2cm}
\setbeamercolor{block title}{bg=dangerred,fg=white}
\begin{block}{Số đoạn mỗi cặp \& length ratio}
\small
\begin{itemize}
  \item[\ding{55}] Trung bình \textbf{30,2} đoạn plagiarism/cặp (median 17, max 1.085) -- phân bố lệch phải mạnh.
  \item[\ding{55}] Length ratio (đoạn nghi vấn / đoạn nguồn) trung vị $\approx$ \textbf{1,08}; obfuscation càng khó, ratio càng phân tán (P90 hard $\approx$ 1,95).
  % \item[\ding{55}] \textbf{similarity} trung bình $\approx$ 0,98 cho cả plagiarism và altered $\rightarrow$ paraphrase LLM rất "khít" nghĩa gốc.
\end{itemize}
\end{block}
\end{frame}
% \begin{frame}{Đánh giá sinh văn bản trên tập Test (249 câu)}
% \begin{table}
% \centering
% \renewcommand{\arraystretch}{1.3}
% \rowcolors{2}{softgray}{}
% \begin{tabular}{lrrrr}
% \toprule
% \textbf{Chỉ số đo lường} & \textbf{Mean} & \textbf{Std} & \textbf{Min} & \textbf{Max}\\
% \midrule
% \textbf{BERTScore F1} & \textbf{0.9304} & 0.0169 & 0.8543 & 0.9749\\
% \textbf{ROUGE-L} & \textbf{0.5842} & 0.0904 & 0.2332 & 0.8330\\
% \bottomrule
% \end{tabular}
% \end{table}

% \vspace{0.2cm}
% \setbeamercolor{block title}{bg=primaryblue,fg=white}
% \begin{block}{Nhận xét chuyên sâu}
% \begin{itemize}
%   \item \textbf{BERTScore F1 (~0.93):} Câu trả lời tạo ra rất tiệm cận về \textit{mặt ngữ nghĩa} so với đáp án gốc (Ground Truth).
%   \item \textbf{ROUGE-L (~0.58):} Độ khớp chuỗi (lexical matching) thấp hơn là dấu hiệu tích cực! LLM diễn đạt uyển chuyển theo văn phong "tư vấn", không học vẹt cứng nhắc văn bản luật nguyên thủy.
%   \item \textbf{Độ trễ thấp (Std):} Phân bố điểm đồng đều chứng minh độ ổn định (consistency) cao của bản LoRA fine-tune.
% \end{itemize}
% \end{block}
% \end{frame}

% \section{Phương pháp tiếp cận}
% {
% \setbeamercolor{background canvas}{bg=white} % Đặt nền trắng

% \begin{frame}[plain]
%     \vfill
%     \centering
%     % Dùng lệnh \textcolor để ép màu xanh hcmusblue cho dòng chữ
%     \textcolor{hcmusblue}{\textbf{\Huge Phương pháp tiếp cận}}
%     \vfill
% \end{frame}
% }



% \centering
% \begin{tikzpicture}[
%   node distance=1.2cm,
%   box/.style={draw=primaryblue!80, fill=primaryblue!10, thick, rounded corners=4pt, align=center, minimum width=4.5cm, minimum height=0.7cm, font=\small\bfseries},
%   io/.style={draw=accentorange!80, fill=accentorange!10, thick, rounded corners=4pt, align=center, minimum width=4cm, minimum height=0.6cm, font=\small},
%   arrow/.style={-Latex, line width=1pt, primaryblue}
% ]
% % Nodes
% \node[io] (q) {Câu hỏi pháp lý};
% \node[box, below of=q] (retrieval) {Hybrid Retrieval\\{\footnotesize BM25 + BGE-M3 + FAISS}};
% \node[box, below of=retrieval] (rrf) {RRF Fusion\\{\footnotesize (Thuật toán hợp nhất thứ hạng)}};
% \node[box, below of=rrf] (rerank) {Cross-Encoder Reranker\\{\footnotesize BAAI/bge-reranker-large}};
% \node[box, below of=rerank] (generator) {Generator Module\\{\footnotesize Qwen2.5-1.5B + LoRA Adapter}};
% \node[io, below of=generator, fill=successgreen!10, draw=successgreen!80] (answer) {Câu trả lời cấu trúc + Căn cứ luật};

% % Arrows
% \draw[arrow] (q) -- (retrieval);
% \draw[arrow] (retrieval) -- node[right, font=\scriptsize] {$Top_K=15$} (rrf);
% \draw[arrow] (rrf) -- (rerank);
% \draw[arrow] (rerank) -- node[right, font=\scriptsize] {$Top_N=3$} (generator);
% \draw[arrow] (generator) -- (answer);
% \end{tikzpicture}
% \end{frame}

% \begin{frame}{Chiến lược Truy xuất (Hybrid Retriever)}
% \begin{columns}[T]
% \begin{column}{0.48\textwidth}
% \setbeamercolor{block title}{bg=primaryblue,fg=white}
% \begin{block}{1. Sparse Retrieval (BM25)}
% \small
% Nắm bắt chính xác các \textbf{từ khóa pháp lý} đặc thù có trong câu hỏi. Sử dụng \textit{rank\_bm25}.
% \end{block}
% \end{column}

% \begin{column}{0.48\textwidth}
% \setbeamercolor{block title}{bg=accentorange,fg=white}
% \begin{block}{2. Dense Retrieval (FAISS)}
% \small
% Truy xuất theo \textbf{tương đồng ngữ nghĩa} sử dụng model \textit{BAAI/bge-m3} (chuyên biệt đa ngôn ngữ).
% \end{block}
% \end{column}
% \end{columns}

% \vspace{0.4cm}
% \textbf{\textcolor{primaryblue}{\rule{\textwidth}{1pt}}}
% \vspace{0.2cm}

% \begin{itemize}
%   \item[\ding{51}] \textbf{Tích hợp RRF (Reciprocal Rank Fusion):} Kết hợp sức mạnh của 2 phương pháp với hệ số mượt $k=60$. Lọc ra $15$ ứng viên tốt nhất.
%   \item[\ding{51}] \textbf{Reranking tinh chỉnh:} Dùng \textit{bge-reranker-large} đánh giá chéo (cross-encoder) giữa câu hỏi và ngữ cảnh, chọn ra \textbf{top-3} điều khoản chính xác nhất nạp vào LLM.
% \end{itemize}
% \end{frame}

% \begin{frame}{Lựa chọn LLM sinh câu trả lời}
% \begin{columns}[T]
%     \begin{column}{0.48\textwidth}
%         \centering
%         \includegraphics[width=\textwidth]{figures/3_model_base.png}
%     \end{column}
%     \begin{column}{0.48\textwidth}
%         \centering
%         \includegraphics[width=\textwidth]{figures/gen_time_3modelbase.png}
%     \end{column}
% \end{columns}

% \vspace{0.3cm}
% \setbeamercolor{block title}{bg=successgreen,fg=white}
% \begin{block}{Insight: Sự cân bằng giữa Chất lượng và Tốc độ}
% \begin{itemize}
%     \item \textbf{Về ngữ nghĩa:} Gemma 2B nhỉnh hơn ở ROUGE-L, nhưng Qwen 1.5B bám rất sát ở BERTScore (0.773 so với 0.761).
%     \item \textbf{Về tốc độ (Yếu tố quyết định):} Qwen 1.5B áp đảo hoàn toàn với thời gian sinh văn bản chỉ \textbf{22.45s}, nhanh hơn Gemma và Llama gấp 1.5 đến 2 lần. 
%     \item \textbf{Kết luận:} Qwen 1.5B là lựa chọn tối ưu nhất: Đủ thông minh để xử lý tiếng Việt và đủ nhẹ để tiết kiệm tài nguyên phần cứng khi huấn luyện LoRA.
% \end{itemize}
% \end{block}
% \end{frame}

% \begin{frame}{Chiến lược Finetune LLM sinh câu trả lời (Generator)}
% \setbeamercolor{block title}{bg=successgreen,fg=white}
% \begin{block}{Tinh chỉnh mô hình (Fine-Tuning)}
% \begin{itemize}
%   \item \textbf{Base Model:} \textit{Qwen/Qwen2.5-1.5B-Instruct}.
%   \item \textbf{Tối ưu hóa:} Tích hợp trọng số \textbf{LoRA Adapter} (chuyên miền Luật Lao động) giúp sinh văn bản chuẩn xác văn phong pháp lý.
% \end{itemize}
% \end{block}

% \vspace{0.2cm}
% \textbf{Kỹ thuật Prompting hướng Luật sư:} \\
% Prompt được thiết kế cứng (System Prompt) ép mô hình trả lời theo nguyên tắc "không bịa luật" với cấu trúc 3 phần rõ rệt:
% \begin{enumerate}
%     \item \textbf{Khẳng định trực tiếp:} Đúng/Sai, Vi phạm/Hợp pháp.
%     \item \textbf{Căn cứ pháp lý:} Trích dẫn trực tiếp Điều/Khoản luật từ ngữ cảnh được cung cấp.
%     \item \textbf{Lời khuyên:} Đưa ra hướng giải quyết thực tiễn cho người hỏi.
% \end{enumerate}
% \end{frame}

% \section{Hệ thống}
% {
% \setbeamercolor{background canvas}{bg=white} % Đặt nền trắng

% \begin{frame}[plain]
%     \vfill
%     \centering
%     % Dùng lệnh \textcolor để ép màu xanh hcmusblue cho dòng chữ
%     \textcolor{hcmusblue}{\textbf{\Huge Hệ thống}}
%     \vfill
% \end{frame}
% }
% \begin{frame}{Triển khai hệ thống thực tế}
% \begin{table}
% \centering
% \renewcommand{\arraystretch}{1.4}
% \begin{tabular}{p{2.5cm} p{8cm}}
% \toprule
% \textbf{\textcolor{primaryblue}{Thành phần}} & \textbf{\textcolor{primaryblue}{Nhiệm vụ \& Công nghệ}}\\
% \midrule
% \texttt{\textbf{backend/}} & Lõi hệ thống viết bằng \textbf{FastAPI}. Đảm nhận RAG engine, tự động nạp cache, Reranker, và Generator (\texttt{main.py}). \\
% \texttt{\textbf{frontend/}} & Giao diện tương tác người dùng thân thiện xây dựng bằng \textbf{Streamlit}: \texttt{app.py}. \\
% \texttt{\textbf{scripts/}} & Pipeline chuẩn bị dữ liệu (\texttt{data\_prep\_agent.py}), chuyển đổi ChatML, sinh vector DB và script đánh giá. \\
% \texttt{\textbf{notebooks/}} & Thử nghiệm Baseline zero-shot cho các model (Qwen, Llama3, Gemma) và huấn luyện LoRA.\\
% \texttt{\textbf{outputs/}} & Lưu trữ LoRA checkpoints (275, 550, 825), adapter, vector database cache, và kết quả test.\\
% \bottomrule
% \end{tabular}
% \end{table}
% \end{frame}

% \section{Đánh giá}
% {
% \setbeamercolor{background canvas}{bg=white} % Đặt nền trắng

% \begin{frame}[plain]
%     \vfill
%     \centering
%     % Dùng lệnh \textcolor để ép màu xanh hcmusblue cho dòng chữ
%     \textcolor{hcmusblue}{\textbf{\Huge Đánh giá}}
%     \vfill
% \end{frame}
% }
% % FRAME 1: ĐÁNH GIÁ RETRIEVER
% \begin{frame}{1. Đánh giá Mô hình Truy xuất (Retriever)}
% \begin{table}
% \centering
% \renewcommand{\arraystretch}{1.3}
% \begin{tabular}{lccc}
% \toprule
% \textbf{Phương pháp truy xuất} & \textbf{Hit@1} & \textbf{Hit@5} & \textbf{MRR@5} \\
% \midrule
% Chỉ dùng BM25 (Sparse) & 0.44 & 0.62 & 0.504 \\
% Chỉ dùng FAISS (Dense) & 0.32 & 0.70 & 0.447 \\
% \textbf{Hybrid (BM25 + FAISS)} & 0.44 & 0.62 & 0.515 \\
% \textbf{Hybrid + Reranker} & \textbf{0.36} & \textbf{0.66} & \textbf{0.484} \\
% \bottomrule
% \end{tabular}
% \end{table}

% \vspace{0.2cm}
% \setbeamercolor{block title}{bg=dangerred,fg=white}
% \begin{block}{Insight: Hiện tượng "Lệch miền dữ liệu" (Domain Shift)}
% Đặc thù văn bản luật yêu cầu khớp từ khóa cực chuẩn, giúp BM25 có Hit@1 rất cao. Khi áp dụng Reranker đa ngữ (Zero-shot), mô hình bị nhầm lẫn giữa các điều luật có ngữ nghĩa na ná nhau, làm tụt Top 1. Điều này chứng minh: \textbf{Cần fine-tune riêng mô hình Reranker trên data Luật VN trong tương lai.}
% \end{block}
% \end{frame}

% FRAME: Evaluation Metrics

\begin{frame}{Đánh giá (1/2)}

\small

\begin{columns}[T]

% Precision
\begin{column}{0.48\textwidth}

\begin{block}{Precision}

\[
\text{Precision}=\frac{TP}{TP+FP}
\]

Đo tỷ lệ ký tự thuộc các vùng đạo văn được dự đoán chính xác trên tổng số ký tự được
hệ thống dự đoán là đạo văn.

\end{block}

\end{column}

% Recall
\begin{column}{0.48\textwidth}

\begin{block}{Recall}

\[
\text{Recall}=\frac{TP}{TP+FN}
\]

Đo khả năng phát hiện các đoạn đạo văn thực sự tồn tại trong tài liệu. Giá trị Recall cao
cho thấy hệ thống ít bỏ sót các trường hợp đạo văn.

\end{block}

\end{column}

\end{columns}

\vspace{-0.2cm}

\begin{columns}[c]

\begin{column}{0.75\textwidth}

\begin{block}{F1-score}

\[
F1=\frac{2\times\text{Precision}\times\text{Recall}}
{\text{Precision}+\text{Recall}}
\]

F1-score là trung bình điều hòa giữa Precision và Recall, phản ánh hiệu quả tổng thể của hệ
thống trong việc cân bằng giữa độ chính xác và khả năng bao phủ.

\end{block}

\end{column}

\end{columns}

\end{frame}


\section{Đánh giá}
\begin{frame}{Đánh giá (2/2)}

\small

\begin{block}{Granularity}

\[
gran(P,D)=
\frac{1}{|P_D|}
\sum_{p\in P_D}|D_p|
\]

Đánh giá mức độ phân mảnh của các dự đoán trong đó $P_D$ là tập các trường hợp đạo văn được phát hiện ít nhất một lần và $D_p$ là số lượng dự đoán tương ứng với trường hợp đạo văn $p$.

$\rightarrow$ Giá trị Granularity lý tưởng bằng 1, tức mỗi trường hợp đạo văn chỉ được phát
hiện bởi một dự đoán.

\end{block}

\vspace{0.3cm}

\begin{block}{PlagDet}

\[
PlagDet(P,D)=
\frac{F1(P,D)}
{\log_2(1+gran(P,D))}
\]

PlagDet là sự kết hợp giữa F1-score và Granularity nhằm đánh giá đồng thời chất lượng phát hiện và mức độ phân mảnh của kết quả. 

$\rightarrow$ PlagDet cao khi hệ thống vừa phát hiện chính xác các đoạn đạo văn vừa hạn chế việc phân mảnh


\end{block}

\end{frame}


% % FRAME 3: HISTOGRAM ĐỘ ỔN ĐỊNH
% \begin{frame}{3. Đánh giá Độ ổn định: Phân phối tần suất RAGAS}
% \centering
% % Hình trên: Histogram Faithfulness
% \includegraphics[height=0.3\textheight, keepaspectratio]{figures/hist_faithfulness.png}

% \vspace{0.1cm}
% % Hình dưới: Histogram Relevancy
% \includegraphics[height=0.3\textheight, keepaspectratio]{figures/hist_relevancy.png}

% \vspace{0.2cm}
% \setbeamercolor{block title}{bg=dangerred,fg=white}
% \begin{block}{Insight: Triệt tiêu các trường hợp sai lệch nghiêm trọng}
% Biểu đồ phân phối tần suất (Histogram) cho thấy rõ phong độ thực tế của mô hình: Bản LoRA đã \textbf{gom tụ dữ liệu} mạnh mẽ về phía dải điểm cao (0.75 - 1.0) và \textbf{xóa sổ hoàn toàn} các cột lỗi ở mốc 0.0 (những câu trả lời hoàn toàn sai luật hoặc lạc đề của mô hình Base).
% \end{block}
% \end{frame}

% FRAME 5: TỔNG QUÁT ƯU NHƯỢC ĐIỂM
% \begin{frame}{Tổng kết: Ưu điểm \& Hạn chế}
% \begin{columns}[T]
% \begin{column}{0.48\textwidth}
% \setbeamercolor{block title}{bg=successgreen,fg=white}
% \begin{block}{Điểm mạnh cốt lõi}
% \begin{itemize}
%   \item[\ding{51}] \textbf{Tối ưu hóa hệ thống:} Áp dụng hiệu quả kỹ thuật fine-tuning LoRA lên mô hình Qwen2.5-1.5B, cho kết quả vượt trội so với Zero-shot.
%   \item[\ding{51}] \textbf{Đánh giá toàn diện:} Kết hợp các độ đo từ vựng học (ROUGE, BERTScore) và đo lường ngữ nghĩa dựa trên LLM (Faithfulness, Relevancy).
%   \item[\ding{51}] \textbf{Tính thực tiễn cao:} Thiết lập prompt sinh câu trả lời theo cấu trúc 3 phần rõ ràng, hỗ trợ tư vấn pháp lý trực tiếp.
% \end{itemize}
% \end{block}
% \end{column}

% \begin{column}{0.48\textwidth}
% \setbeamercolor{block title}{bg=dangerred,fg=white}
% \begin{block}{Hạn chế \& Định hướng}
% \begin{itemize}
%   \item[\ding{55}] \textbf{Quy mô dữ liệu:} Kho ngữ liệu hiện giới hạn ở Luật Lao động, cần mở rộng rà soát thêm các Nghị định và Thông tư hướng dẫn thi hành.
%   \item[\ding{55}] \textbf{Hiệu năng truy xuất:} Giai đoạn chấm điểm lại (Reranking) bằng mô hình Cross-Encoder tiêu tốn nhiều thời gian xử lý (latency).
%   \item[\ding{55}] \textbf{Reranker chưa tối ưu:} Mô hình BGE-Reranker chạy Zero-shot gây hiện tượng lệch miền dữ liệu.
% \end{itemize}
% \end{block}
% \end{column}
% \end{columns}
% \end{frame}

% \section{Kết luận}
% {
% \setbeamercolor{background canvas}{bg=white} % Đặt nền trắng

% \begin{frame}[plain]
%     \vfill
%     \centering
%     % Dùng lệnh \textcolor để ép màu xanh hcmusblue cho dòng chữ
%     \textcolor{hcmusblue}{\textbf{\Huge Kết luận}}
%     \vfill
% \end{frame}
% }

% \begin{frame}{Kết luận \& Hướng triển khai mở rộng}
% \setbeamercolor{block title}{bg=primaryblue,fg=white}
% \begin{block}{Đóng góp của đồ án}
% \small
% Xây dựng thành công hệ thống \textbf{Legal RAG QA System} hoàn chỉnh cho Luật Lao động Việt Nam. Làm chủ toàn bộ quy trình: từ xử lý dữ liệu (ETL), Fine-Tuning mô hình ngôn ngữ lớn (Qwen2.5 + LoRA), tinh chỉnh chiến lược Hybrid Retrieval, cho đến việc triển khai hệ thống phục vụ thực tiễn.
% \end{block}

% \vspace{0.2cm}
% \textbf{\large \textcolor{accentorange}{Các hướng phát triển hệ thống}}
% \vspace{0.1cm}
% \begin{itemize}
%   \item \textbf{Huấn luyện Reranker pháp lý:} Fine-tune mô hình Cross-Encoder trực tiếp trên bộ truy vấn - điều luật VN để triệt tiêu độ lệch miền (Domain Shift).
%   \item \textbf{Mở rộng Tri thức:} Tích hợp các văn bản dưới luật (Nghị định, Thông tư, Án lệ) để bao quát tình huống phức tạp.
%   \item \textbf{Tối ưu suy diễn (Real-time):} Vận hành bằng các framework chuyên dụng (\textit{vLLM} / \textit{TensorRT-LLM}) nhằm giảm triệt để độ trễ khi sinh văn bản.
% \end{itemize}
% \end{frame}

% \begin{frame}{Tài liệu tham khảo}
% \small
% \begin{itemize}
%   \item HCMUS Logo: \url{https://en.hcmus.edu.vn/logo/}
%   \item FIT/HCMUS hướng dẫn viết và trình bày KLTN/TTTN/ĐATN: \url{https://www.fit.hcmus.edu.vn/vn/Default.aspx?id=13153&tabid=292}
%   \item FIT/HCMUS hướng dẫn bảo vệ: \url{https://www.fit.hcmus.edu.vn/vn/UserFiles/7562_FIT-HCMUS_Huong-dan-bv-KLTN-TTTN-TTDATN-K2020-1.pdf}
% \end{itemize}
% \end{frame}

\section{Kế hoạch triển khai}
{
\setbeamercolor{background canvas}{bg=white}

\begin{frame}[plain]
    \vfill
    \centering
    \textcolor{hcmusblue}{\textbf{\Huge Kế hoạch triển khai}}
    \vfill
\end{frame}
}


\begin{frame}{Kế hoạch triển khai}
\centering

\begin{tikzpicture}[
    % Sử dụng thư viện positioning để căn khoảng cách tự động giữa các mép hộp
    node distance=0.25cm and 0.5cm, 
    every node/.style={font=\small, align=center},
    % Style cho các Phase thông thường
    phase/.style={
        draw=primaryblue!70,
        fill=primaryblue!5,
        rounded corners=6pt,
        minimum width=8.5cm,
        minimum height=0.9cm,
        line width=0.8pt,
        text width=8.2cm,
        inner sep=4pt
    },
    % Style nổi bật cho Milestone/Phase cuối
    milestone/.style={
        draw=successgreen!70,
        fill=successgreen!5,
        rounded corners=6pt,
        minimum width=8.5cm,
        minimum height=0.9cm,
        line width=1pt,
        text width=8.2cm,
        inner sep=4pt
    },
    % Mũi tên thanh lịch hơn
    arrow/.style={
        -{Stealth[scale=1.2]},
        line width=1.2pt,
        primaryblue!80
    }
]

% Nodes - Sắp xếp đúng thứ tự từ trên xuống dưới để tránh lỗi "below=of"
\node[phase] (p1) {\textbf{Phase 1: Dataset Preparation \& Data Pipeline}};

\node[phase, below=of p1] (p2) 
{\textbf{Phase 2: Knowledge Base}};

\node[phase, below=of p2] (p3) 
{\textbf{Phase 3: Source Retrieval} \\ \scriptsize\textcolor{gray}{(BM25 + Dense Retrieval + FAISS)}};

\node[phase, below=of p3] (p4) 
{\textbf{Phase 4: Text Alignment} \\ \scriptsize\textcolor{gray}{(Detect plagiarism spans)}};

\node[phase, below=of p4] (p5) 
{\textbf{Phase 5: Plagiarism Scoring} \\ \scriptsize\textcolor{gray}{(0--1 score, plagiarism percentage, source identification)}};

\node[milestone, below=of p5] (p6) 
{\textbf{Phase 6: Evaluation \& Optimization} \\ \scriptsize\textcolor{successgreen!80!black}{(Precision, Recall, F1, Granularity, PlagDet)}};

% Arrows (Vẽ các mũi tên nối kết)
\draw[arrow] (p1) -- (p2);
\draw[arrow] (p2) -- (p3);
\draw[arrow] (p3) -- (p4);
\draw[arrow] (p4) -- (p5);
\draw[arrow] (p5) -- (p6);

\end{tikzpicture}

\vspace{0.2cm}

% Thiết kế Block Milestones gọn gàng
\setbeamercolor{block title}{bg=primaryblue,fg=white}
\setbeamercolor{block body}{bg=primaryblue!5,fg=black}
\begin{block}{\rule{0pt}{2.0ex}✔ Trạng thái thực hiện (Milestones)}
\small
\vspace{0.1cm}
\begin{columns}[T]
    \begin{column}{0.5\textwidth}
        \begin{itemize}
            \item[\textcolor{successgreen}{$\checkmark$}] \textbf{P1:} Preprocessing \& Pipeline.
            \item[\textcolor{successgreen}{$\checkmark$}] \textbf{P2:} Retrieval baseline.
            \item[\textcolor{successgreen}{$\checkmark$}] \textbf{P3:} Alignment (Suspicious vs. Source).
        \end{itemize}
    \end{column}
    \begin{column}{0.5\textwidth}
        \begin{itemize}
            \item[\textcolor{successgreen}{$\checkmark$}] \textbf{P4:} Tính điểm \& Truy vết.
            \item[\textcolor{successgreen}{$\checkmark$}] \textbf{P5:} Đánh giá PlagDet \& Tối ưu.
        \end{itemize}
    \end{column}
\end{columns}
\vspace{0.1cm}
\end{block}

\end{frame}
{
\setbeamercolor{background canvas}{bg=white} % Đặt nền trắng

\begin{frame}[plain]
    \vfill
    \centering
    % Dùng lệnh \textcolor để ép màu xanh hcmusblue cho dòng chữ
    \textcolor{hcmusblue}{\textbf{\Huge Cảm ơn thầy và các bạn đã lắng nghe}}
    \vfill
\end{frame}
}


\end{document}
